% answers_short.tex —— 1977 年·贵州卷 · 简短答案速查
% 由 tools/make_short_answers.py 自 answers.tex 的 \gksolhead/\gkanswer 提取，
% 并与 1977/1977_数学真题_贵州_答案要点.md 逐题比对后生成，勿手改（改源 MD 后重跑）。

\begin{gknotice}
\textbf{使用说明}：本节只给最终答案便于快速核对，完整推导见解析版。答案为编写者独立推导，\textbf{非官方标准答案}。
\end{gknotice}

\gkpart{题目}
\begin{itemize}[leftmargin=2.6em]
  \item[\textbf{1.}] (1) $b^{\frac{5}{12}}$（即 $\sqrt[12]{b^{5}}$）；(2) $\dfrac{1}{2}$
  \item[\textbf{2.}] (1) $\sin\alpha=\pm\dfrac{2ab}{a^{2}+b^{2}}$；(2) $m^{2}+\dfrac{1}{m^{2}}=2$，$m^{3}+\dfrac{1}{m^{3}}=2$
  \item[\textbf{3.}] 见解析版
  \item[\textbf{4.}] $3x-y-6=0$（即 $y=3x-6$）。
  \item[\textbf{5.}] 当长、宽相等且均为 $2$ 米（边长 $2$ 米的正方形）时，窗子透过的光线最多。
  \item[\textbf{6.}] $a\cot\dfrac{\alpha}{2}$。
  \item[\textbf{7.}] $q\neq1$ 时原式等于 $q$；$q=1$ 时原式等于 $1$。
  \item[\textbf{8.}] $a>0$ 时（L）与（c）相交；$a=0$ 时相切；$a<0$ 时不相交。
\end{itemize}
